\section{C7 Relation-Resistant Compaction and Final Validation Status}
\label{app:c7-relation-resistant}

C7 is the final construction pass in this validation package.  It responds to
C6's negative finding: complete projective C4 blocks are correct, but they
publish many weight-three public relations.  C7 therefore changes the main
candidate compactor from a full projective block dictionary to a
relation-resistant coordinate compactor, while keeping screened random masks as
an experimental optimization.

\subsection{C7 coordinate compactor}

Partition the final GSW last-row coefficient vector into blocks
$\alpha_B\in\F_q^b$.  In the C7 coordinate mode, the public basis for a block
is simply
\[
    U_B = \{e_1,\ldots,e_b\},
\]
where $e_i$ are the standard coordinate vectors.  The public key stores
\[
    \CLPN.\Enc_r(\langle e_i,s_B\rangle)
    =
    \CLPN.\Enc_r(s_{B,i})
\]
for each coordinate in the block.  Evaluation decomposes
\[
    \alpha_B=\sum_{i=1}^b \alpha_i e_i
\]
and linearly combines the corresponding CLPN ciphertexts.

This mode is intentionally conservative.  It does not reduce the worst-case
number of compaction terms below the legacy coordinate compactor, but it also
does not create the complete projective-line relation surface from C4.  In
particular, inside a block the matrix whose rows are $e_1,\ldots,e_b$ has full
row rank and no non-zero kernel relation.

\begin{lemma}[C7 coordinate coverage]
For every $\alpha_B\in\F_q^b$, the C7 coordinate basis gives a decomposition
of Hamming weight exactly $\wt(\alpha_B)$ and at most $b$.
\end{lemma}

\begin{proof}
The coordinate identity
$\alpha_B=\sum_i \alpha_i e_i$ is unique.  Zero coefficients contribute no
term, so the number of non-zero terms is $\wt(\alpha_B)\leq b$.
\end{proof}

\begin{lemma}[No within-block kernel relation]
The C7 coordinate basis has no non-zero linear relation among its public masks.
\end{lemma}

\begin{proof}
If $\sum_i z_i e_i=0$, then the $i$th coordinate of the sum is $z_i$.  Hence
all $z_i=0$.
\end{proof}

\subsection{Experimental screened-random mode}

C7 also includes an optional screened-random mode.  Each block begins with the
coordinate basis so that correctness is guaranteed.  The generator then tries
to add dense random masks.  A candidate mask is rejected if it is projectively
duplicate or if adding it creates a local kernel relation up to the configured
screen weight.  During compaction, the evaluator first searches for a sparse
representation using the extra masks and falls back to the coordinate
representation when no sparse representation is found.

This mode is not the main security claim.  It is an optimization target for
future work because every added mask increases the public LPN sample surface.
The final package therefore marks C7 coordinate as the main research default
and screened-random C7 as experimental.

\subsection{Correctness theorem}

Let $C_\star$ be the evaluated GSW ciphertext and let $r_\star$ be its last
row.  Let $B_1,\ldots,B_L$ be the coefficient blocks touched by $r_\star$.
C7 coordinate compaction produces
\[
    \sum_{B}\sum_{i\in B} r_{\star,i}\,\CLPN.\Enc_r(s_i),
\]
which is a CLPN encryption of
\[
    \sum_i r_{\star,i}s_i = r_\star\cdot \widetilde{s}.
\]
As in the main construction, if
\[
    r_\star\cdot\widetilde{s}=f(\mu_1,\ldots,\mu_m)+e_\star,
\]
then the compacted ciphertext decrypts to the correct value whenever
$e_\star=0$ and the CLPN decoder succeeds.

If $T$ CLPN ciphertexts are added during C7 compaction, the q-ary symmetric
error rate entering the repetition decoder is bounded by
\[
    \eta_T
    =
    \frac{q-1}{q}
    \left(
    1-
    \left(1-\frac{q\eta_c}{q-1}\right)^T
    \right).
\]
For coordinate C7, $T\leq \wt(r_\star)\leq N$.  This is worse than one-term
projective C4 in pure noise terms, but it removes the complete projective
weight-three relation surface.

\subsection{Validation results in the repository}

The v0.9-C7 package contains the following additional validation artifacts:
\begin{itemize}
    \item \texttt{src/sable/c7\_relation\_resistant.py}: C7 coordinate and screened-random compaction;
    \item \texttt{experiments/run\_c7\_arithmetic\_suite.py}: all arithmetic and Boolean-operation tests under C7;
    \item \texttt{experiments/run\_c7\_compare\_baselines.py}: operation-count comparison against TFHE/FHEW, BFV/BGV, and CKKS-style families;
    \item \texttt{experiments/run\_c7\_readiness.py}: final readiness report;
    \item \texttt{docs/c7\_readiness\_output.json}: machine-readable status summary.
\end{itemize}

The current C7 validation run reports that all tested arithmetic operations
succeed in the toy-clean configuration: addition, subtraction, negation,
public scalar multiplication, public constants, multiplication, squaring,
affine combinations, dot products, polynomial evaluation, balanced products,
quadratic forms, and the Boolean gates NOT, AND, OR, XOR, NAND, NOR, XNOR, and
implication encoded over $\F_q$.  The test suite reports \texttt{95 passed}.

\subsection{Final project status}

At the end of C7, the project is ready to pause as a research prototype and
manuscript package.  The main candidate is no longer C4 projective compaction;
it is C7 coordinate relation-resistant compaction.  The package is not ready
for deployment or standardization, and it does not claim certified concrete
security.  The remaining scientific work is independent cryptanalysis,
stronger q-ary decoding beyond repetition coding, and measured benchmarks
against optimized HE implementations on matched circuits.
